\documentclass[12pt]{jlreq}
\usepackage{amsmath, amssymb}

\newcommand{\C}{\mathbb{C}}
\newcommand{\R}{\mathbb{R}}
\newcommand{\Z}{\mathbb{Z}}
\newcommand{\N}{\mathbb{N}}
\newcommand{\F}{\mathbb{F}}
\newcommand{\Prob}{\mathbb{P}}
\newcommand{\Exp}{\mathbb{E}}
\newcommand{\dd}{\,d}
\newcommand{\Ker}{\operatorname{Ker}}
\newcommand{\Impart}{\operatorname{Im}}
\newcommand{\Hom}{\operatorname{Hom}}

\begin{document}

\(\#A\) で集合 \(A\) の元の個数を表す。群 \(G\) とその部分群 \(H\) および \(L\) に対して，次のように定める。
\[
  HgL=\{hg\ell\mid h\in H,\ \ell\in L\}\quad (g\in G),
\]
\[
  H\backslash G/L=\{HgL\mid g\in G\}.
\]
\(2\) 以上の正整数 \(n\) に対して \(n\) 次対称群 \(\mathfrak S_n\) を考える。さらに整数 \(1\le k\le n\) に対して，
\[
  H_k=\{\sigma\in\mathfrak S_n\mid \sigma(i)\in\{1,\ldots,k\}\ (1\le i\le k)\}
\]
と定めると，\(H_k\) は \(\mathfrak S_n\) の部分群になる。以下の問に答えよ。

\begin{enumerate}
  \item \(\max\{\#H_1gH_1\mid g\in\mathfrak S_n\}\) を求めよ。
  \item \(1\le p,q\le n\) なる整数 \(p,q\) に対して \(\#H_p\backslash\mathfrak S_n/H_q\) を求めよ。
\end{enumerate}

\end{document}